\documentclass[11pt,a4paper]{article}
\usepackage[utf8]{inputenc}
\usepackage[T1]{fontenc}
\usepackage[brazil]{babel}
\usepackage{amsmath,booktabs,longtable}
\usepackage[a4paper,margin=2cm]{geometry}
\usepackage[backend=biber,style=authoryear]{biblatex}
\addbibresource{referencias.bib}
\title{Ações e fundos imobiliários: patrimônio e renda em 10 e 20 anos}
\author{Maycon Douglas Alves Moraes}
\date{30 de outubro de 2026}
\begin{document}
\maketitle
\begin{abstract}
Compara-se um aporte hipotético de R$ 10.000,00 em quinze ações e quinze fundos imobiliários. Para dez anos, calculam-se patrimônio final e renda potencial pela aplicação do DY recente. O horizonte de vinte anos é discutido, mas não recebe valores inventados: as fontes consultadas não oferecem retorno total homogêneo e auditável para os quinze papéis, e não há quinze FIIs com vinte anos completos de histórico individual comparável.
\end{abstract}
\section{Metodologia}
Para retorno total acumulado R_n, o patrimônio final é C_n=10.000(1+R_n). A taxa anual composta é CAGR=(C_n/10.000)^{1/n}-1. A renda estimada é D_ano=C_n DY_12m e D_mês=D_ano/12. O DY é retrospectivo, não uma garantia. A renda calculada aplica o yield atual ao patrimônio final e não representa os proventos efetivamente recebidos durante a década.
As séries de retorno de ações foram extraídas principalmente do Meus Dividendos, com complementos do Dados B3. O Dados B3 informa combinar cotações da B3 e proventos, com ajustes para eventos corporativos \parencite{dadosb3}. As páginas do Meus Dividendos apresentam a série de cotação e proventos por ativo \parencite{meusdividendos}. Os DY de ações usam principalmente o painel Leitura Singular, referência 29/10/2026. Para FIIs, retornos e rendimentos foram extraídos das páginas individuais do Meus Dividendos, consultadas em 30/10/2026.
\section{Ações: simulação de dez anos}
Aporte de R$ 10.000,00 em aproximadamente outubro de 2016. Valores em reais, arredondados.
\begin{longtable}{lrrrrr}
\toprule Código & Retorno & Patrimônio & DY & Renda anual & Renda mensal\\
\midrule
PETR4 & 614,83\% & 71.483 & 6,89\% & 4.925 & 410\\
VALE3 & 578,73\% & 67.873 & 9,14\% & 6.204 & 517\\
ITUB4 & 193,03\% & 29.303 & 7,96\% & 2.333 & 194\\
BBDC4 & 67,81\% & 16.781 & 11,00\% & 1.846 & 154\\
CMIG4 & 253,04\% & 35.304 & 11,24\% & 3.968 & 331\\
TAEE11 & 217,51\% & 31.751 & 7,57\% & 2.404 & 200\\
EGIE3 & 115,40\% & 21.540 & 4,10\% & 883 & 74\\
CPFE3 & 151,34\% & 25.134 & 8,69\% & 2.184 & 182\\
WEGE3 & 794,90\% & 89.490 & 4,24\% & 3.794 & 316\\
ABEV3 & 8,12\% & 10.812 & 5,56\% & 601 & 50\\
ITSA4 & 250,40\% & 35.040 & 8,86\% & 3.105 & 259\\
BBSE3 & 130,56\% & 23.056 & 10,96\% & 2.527 & 211\\
VIVT3 & 115,34\% & 21.534 & 4,29\% & 924 & 77\\
GGBR4 & 401,83\% & 50.183 & 2,49\% & 1.250 & 104\\
PSSA3 & 309,28\% & 40.928 & 6,42\% & 2.628 & 219\\
\bottomrule
\end{longtable}
\section{FIIs: simulação de dez anos}
Doze fundos têm retorno e DY consolidados nesta versão. HGRE11, HGCR11 e SDIL11 ficam pendentes de validação.
\begin{longtable}{lrrrrr}
\toprule Código & Retorno & Patrimônio & DY & Renda anual & Renda mensal\\
\midrule
HGLG11 & 140,95\% & 24.095 & 8,20\% & 1.976 & 165\\
KNCR11 & 78,03\% & 17.803 & 15,10\% & 2.688 & 224\\
KNRI11 & 66,45\% & 16.645 & 7,90\% & 1.315 & 110\\
MXRF11 & 115,69\% & 21.569 & 13,20\% & 2.847 & 237\\
HGBS11 & 62,99\% & 16.299 & 9,62\% & 1.568 & 131\\
BRCR11 & 13,26\% & 11.326 & 5,90\% & 668 & 56\\
JSRE11 & 23,80\% & 12.380 & 5,80\% & 718 & 60\\
RNGO11 & 41,15\% & 14.115 & 8,84\% & 1.248 & 104\\
FIIB11 & 178,09\% & 27.809 & 8,20\% & 1.502 & 125\\
HTMX11 & 97,35\% & 19.735 & 11,26\% & 2.222 & 185\\
BCRI11 & 47,09\% & 14.709 & 9,76\% & 1.436 & 120\\
KNIP11 & 68,84\% & 16.884 & 11,00\% & 1.857 & 155\\
\bottomrule
\end{longtable}
\section{Horizonte de vinte anos}
A janela 30/10/2006--29/10/2026 exige séries contínuas de preços, proventos, eventos corporativos, amortizações e liquidações. As fontes consultadas não disponibilizaram retorno total uniforme para os quinze papéis. Para FIIs, diversos fundos sequer existiam em 2006. Misturar retorno total de dez anos com retorno apenas de preço em vinte anos seria inconsistente. Portanto, o resultado correto nesta versão é: não calculável com dados comparáveis validados. Extrapolar o CAGR de dez anos seria um cenário hipotético, não resultado histórico.
\section{Interpretação e limitações}
Entre as ações com dados consolidados, WEGE3 apresenta o maior patrimônio hipotético, seguida por PETR4 e VALE3. Entre os FIIs com números, HGLG11, MXRF11 e FIIB11 apresentam os maiores patrimônios finais. O DY atual pode alterar a ordem da renda potencial: o ativo que mais valorizou não precisa ser o que mais renda gera. A amostra é retrospectiva, sujeita a viés de seleção e sobrevivência. Valores nominais, sem custos e tributos; DY não garante renda futura.
\section{Conclusão}
A janela de dez anos permite comparar patrimônio e renda potencial com os dados disponíveis. A janela de vinte anos não foi preenchida porque as séries comparáveis necessárias não foram validadas para quinze ativos de cada classe. Essa limitação é preferível a publicar extrapolações como resultados históricos.
\printbibliography
\end{document}